\honest{A Prodigy of Refutation}{Eliezer Yudkowsky}{2008}

In 1996 three things carried me into my mistake: a death spiral around ``intelligence,'' technophilia, and science fiction that taught me to draw my in-group broadly. Hence my email address, sentience@pobox.com. \nb{The young Yudkowsky's posts to the Extropians mailing list are signed with it from late November 1996.}

Eliezer1996 set out to build superintelligence for the good of humanity and all sentient life. At first, whether it would be good did not occur to me as a separate question. Surely no supermind would be stupid enough to turn the galaxy into paperclips, and surely it would know what is right ``far better than a human being could.'' \nb{The archive agrees. On 22 November 1996 the young Yudkowsky wrote: ``The Powers will be ethical.''}

On a transhumanist mailing list people objected, as I remember it: morality is arbitrary, so a superintelligence would form its own; everyone looks after their own interest, so a superintelligence would seize all the resources; I am human, so I favor humans, and we should upload humans instead of building AI; no one should build an AI without a control system that stops it doing anything bad. Eliezer1996 thought all of this ``obviously wrong'' and kicked the arguments to pieces. I have mostly done this in other posts, and leave the rest ``as an exercise to the reader.'' \nb{The archive has objections of the last two kinds, and the replies: ``THIS LINE OF THOUGHT IS THE GREATEST KNOWN DANGER TO HUMANITY!'' (8 December 1996). Two of the conclusions are close to Yudkowsky's own view in 2008: that no argument moves ``all physically possible minds,'' and that an unfriendly AI would see us as ``atoms which it can use for something else.''}

Why did winning feel like being right? Eliezer1996 was a Traditional Rationalist, who saw science as a fair fight won by the best arguments. He did not reason that if the world's stupidest man says the sun is shining, it must be dark. But finding flaws is easier than finding the truth, and Eliezer1996 was very good at finding flaws. So am I, and refusing to care about flaws would not remove the danger. His side seemed to be winning, so why switch? I quote my ``Twelve Virtues of Rationality'': beginners ``win arguments and acquire an exaggerated view of their own abilities. But it is useless to be superior: Life is not graded on a curve. The best physicist in ancient Greece could not calculate the path of a falling apple.'' Adequacy may be out of reach even with your hardest effort, so spare no thought for whether others are doing worse.

So you cannot rely on anyone else to argue you out of your mistakes or to save you. You alone must find the flaws in your positions, and if you put that burden down, nobody else will pick it up. I wonder if that advice helps anyone before they have blown off their own foot, thinking all the while, correctly, that they were winning the argument. Today I take Nature, not any human being, as my opponent, because human opponents lead to overconfidence. Nature does not match Her problems to your skill, and does not offer a fair chance in return for diligent effort.

No opponent gave Eliezer1996 ``the full reduction of morality,'' a step-by-step account of the algorithms that make morality feel like a fact. That nothing less would have persuaded measures Eliezer1996's level. The few philosophers present could not say that morality was settled, since it is still an open question in philosophy. So they offered classic arguments on all sides, which Eliezer1996 knocked down, and so concluded that philosophy is a wasteland.

It gets worse. Perhaps in 1997, the younger me set out to argue ``inescapably'' that creating superintelligence is right. To be continued.
